Gaussian noise sigma = 0.2; the DAE (same architecture as AE-1024, noisy input $\rightarrow$ clean target, noise redrawn every batch, no labels) is compared using the fixed Q1(a) CNN. Test accuracy: clean 87.6\%, corrupted 17.2\%, blur 30.2\%, clean-AE reconstruction 44.2\%, DAE reconstruction 41.7\% (ceiling: clean AE on clean images 59.9\%). The DAE recovers F = 0.35 of the accuracy lost to noise (95\% CI of the gain +23.3 pp to +25.7 pp), and 0.57 of what is recoverable through this reconstruction architecture. Against the clean AE it gains -2.5 pp, and against blur +11.4 pp. Decision: denoising PARTIALLY RECOVERS the information the classifier needs. Of the remaining 46.0 pp shortfall from clean accuracy, 27.7 pp also occurs when the clean AE reconstructs clean images, so most of it is the classifier's penalty for smooth reconstructions, not noise left in the image.
